% section: この章の地図

\section*{この章の地図}

第2章では、漸化式をいくつかの型に分け、それぞれを計算して解いた。
第3章では、問題を見たときに、差を取るのか、比を取るのか、不動点を引くのか、
あるいは別の表現へ移るのかを考えた。

ここからは、これらの解法を少し遠くから眺める。

漸化式は、単に数列の項を計算するための式ではない。隣り合う項の差を調べれば、
離散的な微分と積分が現れる。同じ関数を繰り返し作用させれば、反復、不動点、
周期、安定性が現れる。線形な漸化式をまとめれば、行列や固有値が現れる。
無限個の項を一つの式にまとめれば、母関数が現れる。そして、数列の規則を
より広い変数へ延長しようとすると、特殊関数が現れる。

この章の目的は、これらの理論をすべて完成させることではない。
第3章までで使った解法が、どのような数学の考え方につながっているのかを
見えるようにすることである。

\subsection*{この章で見る四つの視点}

\paragraph{変化を見る}

\[
  \Delta a_n=a_{n+1}-a_n
\]

差分を取ると、数列がどのように変化しているかが見える。
階差型漸化式を解くことは、差分を元に戻すことでもある。

\paragraph{反復を見る}

\[
  a_{n+1}=f(a_n)
\]

という形は、関数 $f$ を同じ点に何度も作用させる操作である。
一般項が求まらなくても、不動点へ近づくか、周期運動になるか、発散するかを調べられる。

\paragraph{構造を見る}

線形漸化式では、数列をベクトルや行列でまとめることができる。
母関数を使えば、無限個の項の関係を一つの式に移せる。

\paragraph{表現を広げる}

三角関数、階乗、積分、直交多項式などは、漸化式に現れる規則を別の形で保存したものとして見ることができる。

以降では、各節を独立した話題として暗記するのではなく、同じ漸化式を異なる言葉へ
翻訳していく流れとして読む。
